%-------------------------
% Resume in Latex
% Author : Karnayana Rohith
% License : MIT
%------------------------

\documentclass[letterpaper,10.5pt]{article}

\usepackage{latexsym}
\usepackage[empty]{fullpage}
\usepackage{titlesec}
\usepackage{marvosym}
\usepackage[usenames,dvipsnames]{color}
\usepackage{verbatim}
\usepackage{enumitem}
\usepackage[hidelinks]{hyperref}
\usepackage{fancyhdr}
\usepackage[english]{babel}
\usepackage{tabularx}
\usepackage{geometry}

\geometry{
  letterpaper,
  left=0.5in,
  right=0.5in,
  top=0.45in,
  bottom=0.45in
}

\pagestyle{fancy}
\fancyhf{} % clear all header and footer fields
\fancyfoot{}
\renewcommand{\headrulewidth}{0pt}
\renewcommand{\footrulewidth}{0pt}

% Use standard margins
\urlstyle{same}

\raggedbottom
\raggedright
\setlength{\tabcolsep}{0in}

% Sections formatting
\titleformat{\section}{
  \vspace{-10pt}\scshape\raggedright\large\bfseries
}{}{0em}{}[\color{black}\rule{\textwidth}{0.5pt}\vspace{-6pt}]

%-------------------------
% Custom commands
\newcommand{\resumeItem}[1]{
  \item\small{
    {#1 \vspace{-2.5pt}}
  }
}

\newcommand{\resumeSubheading}[4]{
  \vspace{-2pt}\item
    \begin{tabular*}{0.97\textwidth}[t]{l@{\extracolsep{\fill}}r}
      \textbf{#1} & #2 \\
      \textit{\small#3} & \textit{\small #4} \\
    \end{tabular*}\vspace{-6pt}
}

\newcommand{\resumeProjectHeading}[2]{
    \item
    \begin{tabular*}{0.97\textwidth}{l@{\extracolsep{\fill}}r}
      \small#1 & #2 \\
    \end{tabular*}\vspace{-6pt}
}

\newcommand{\resumeSubItem}[1]{\resumeItem{#1}\vspace{-4pt}}

\renewcommand\labelitemii{$\vcenter{\hbox{\tiny$\bullet$}}$}

\newcommand{\resumeSubHeadingListStart}{\begin{itemize}[leftmargin=0.15in, label={}]}
\newcommand{\resumeSubHeadingListEnd}{\end{itemize}}
\newcommand{\resumeItemListStart}{\begin{itemize}}
\newcommand{\resumeItemListEnd}{\end{itemize}\vspace{-5pt}}

%-------------------------------------------
%%%%%%  RESUME STARTS HERE  %%%%%%%%%%%%%%%%############


\begin{document}

%----------HEADING----------
\begin{center}
    \textbf{\Huge \scshape Karnayana Rohith} \\ \vspace{4pt}
    \small Visakhapatnam, Andhra Pradesh, India \\ \vspace{2pt}
    \small +91 97011 28085 $|$ \href{mailto:karnayanarohith@gmail.com}{\underline{karnayanarohith@gmail.com}} $|$ 
    \href{https://linkedin.com/in/rohith-karnayana}{\underline{linkedin.com/in/rohith-karnayana}} $|$ 
    \href{https://github.com/karnayanarohith}{\underline{github.com/karnayanarohith}} $|$ 
    \href{https://rohithkarnayana.vercel.app}{\underline{rohithkarnayana.vercel.app}}
\end{center}


%-----------PROFILE-----------
\section{Profile}
\small{
Final-year B.Tech ECE student specializing in low-level systems security and agentic AI. Experienced in firmware exploitation, kernel debugging, and orchestrating multi-agent LLM systems. Targeting SOC L1 / Cybersecurity Analyst roles to apply hands-on skills in threat mitigation, log analysis, and secure system design.
}


%-----------EDUCATION-----------
\section{Education}
  \resumeSubHeadingListStart
    \resumeSubheading
      {Andhra University College of Engineering}{Visakhapatnam, Andhra Pradesh}
      {Bachelor of Technology (B.Tech) in Electronics \& Communication Engineering $|$ CGPA: 8.4/10}{2023 -- 2027}
    \resumeSubheading
      {Sri Gayathri Junior College}{Gotlam, Vizianagaram, Andhra Pradesh}
      {Intermediate (MPC) $|$ Percentage: 95\%}{2021 -- 2023}
    \resumeSubheading
      {ZPHS Siripuram}{Siripuram, Srikakulam, Andhra Pradesh}
      {Secondary School Certificate (SSC) $|$ Percentage: 85\%}{2021}
  \resumeSubHeadingListEnd


%-----------TECHNICAL SKILLS-----------
\section{Technical Skills}
 \begin{itemize}[leftmargin=0.15in, label={}]
    \small{\item{
     \textbf{Languages}{: C, Python, Bash, Assembly, JavaScript, SQL (PostgreSQL)} \\
     \textbf{Cybersecurity}{: CVE Research \& Exploitation, Log Analysis \& Forensics, MITRE ATT\&CK, Threat Triage \& Incident Response, IAM} \\
     \textbf{AI / ML}{: LangGraph (Agentic AI), Ollama (Local LLMs), Vector DBs (ChromaDB), FastAPI, PyTorch, OpenCV} \\
     \textbf{Systems \& OS}{: Embedded Linux, Android Bootloader Internals, process/memory debugging, filesystem structures} \\
     \textbf{Tools \& Utilities}{: Git, Docker, Linux CLI, GDB, ADB, Fastboot, Wireshark, Splunk, Wazuh}
    }}
 \end{itemize}


%-----------PROJECTS-----------
\section{Projects}
    \resumeSubHeadingListStart
      \resumeProjectHeading
          {\textbf{Aegis} $|$ \emph{Python, LangGraph, Ollama, ChromaDB, Docker}}{}
          \resumeItemListStart
            \resumeItem{Developed an autonomous, dual-mode Blue/Red Team security agent orchestrated with LangGraph and local LLMs via Ollama.}
            \resumeItem{Engineered a RAG-based memory pipeline using ChromaDB to store and recall threat intelligence, vulnerability logs, and attack patterns.}
            \resumeItem{Built and containerized modular tools for automated network reconnaissance, log analysis, and incident triage mapped to the MITRE ATT\&CK framework.}
          \resumeItemListEnd
          
      \resumeProjectHeading
          {\textbf{EXT4 Forensic Recovery Tool} $|$ \emph{C, Linux, Systems Programming}}{}
          \resumeItemListStart
            \resumeItem{Developed a low-level forensic recovery tool in C that parses raw EXT4 disk partitions directly without mounting the filesystem.}
            \resumeItem{Programmed manual superblock, inode table, and directory entry parsers to locate and reconstruct unlinked data block chains.}
            \resumeItem{Designed an ncurses-based TUI with a read-only preview engine to ensure forensically sound, non-destructive file recovery operations.}
          \resumeItemListEnd
    \resumeSubHeadingListEnd


%-----------EXPERIENCE-----------
\section{Experience}
  \resumeSubHeadingListStart

    \resumeSubheading
      {Android Security \& Systems Research}{May 2026 -- July 2026}
      {Self-Directed Research, Realme C15 (MT6765)}{}
      \resumeItemListStart
        \resumeItem{\textbf{CVE \& Vulnerability Profiling:} Conducted MediaTek BROM exploit research to bypass Download Agent authentication using MTKClient and dump the device's seccfg partition. Mapped vulnerability attack surfaces on EOL firmware, profiling CVE-2022-20421 (Binder UAF) and CVE-2024-20106 (Type Confusion) via GDB and dmesg.}
        \resumeItem{\textbf{Custom ROM \& Recovery Deployment:} Executed bootloader unlocking via BROM RAM patching and resolved Android dynamic partition allocation errors during custom recovery flashing. Deployed LineageOS and set up a Kali NetHunter chroot environment to convert the device into a mobile penetration testing platform.}
      \resumeItemListEnd

    \resumeSubheading
      {Network Operations Lead}{March 2026}
      {ELEXCENTRA Hackathon, Andhra University College of Engineering}{}
      \resumeItemListStart
        \resumeItem{Engineered and managed the network infrastructure for a live, 5-hour hackathon supporting 400+ concurrent devices across 6 rooms.}
        \resumeItem{Configured and deployed Wi-Fi and physical LAN cabling using TP-Link switches along with D-Link and Netgear routers.}
        \resumeItem{Diagnosed and resolved real-time network congestion, signal degradation, and routing conflicts under pressure to ensure zero downtime.}
      \resumeItemListEnd

    \resumeSubheading
      {Machine Learning Research Assistant}{2025}
      {NEXUS Research Project, Andhra University (under Dr. Karri Chiranjeevi)}{}
      \resumeItemListStart
        \resumeItem{Engineered a real-time small object detection pipeline for industrial seals and tags utilizing PyTorch, CUDA, and the MMDetection framework.}
        \resumeItem{Benchmarked advanced CNN and Vision Transformer architectures including TOOD, Deformable DETR, and DINO.}
        \resumeItem{Achieved ~98\% accuracy on large objects and improved small-object mAP from 0.0 to 0.28 via resolution scaling and custom data balancing.}
      \resumeItemListEnd

  \resumeSubHeadingListEnd


%-----------ACHIEVEMENTS & HONORS-----------
\section{Achievements \& Honors}
  \resumeItemListStart
    \resumeItem{\textbf{1st Place Champion}, HackAP (2025): Co-engineered 'Foundra' (\href{https://github.com/karnayanarohith/Foundra}{\underline{github.com/karnayanarohith/Foundra}}), an AI startup advisor orchestrating 11 autonomous LLM agents across 6 business development phases.}
    \resumeItem{\textbf{Security Solver} (picoCTF, TryHackMe, OverTheWire): Actively practicing hands-on exploitation and offensive security. Solving 'THE GAUNTLET' (\href{https://github.com/karnayanarohith/THE-GAUNTLET}{\underline{github.com/karnayanarohith/THE-GAUNTLET}}) --- a 21-challenge security lab with public write-ups.}
    \resumeItem{\textbf{Project Finalist}, Smart India Hackathon \& AURORA 2.0 (2024): Developed computer vision models to analyze satellite imagery for illegal mining detection.}
  \resumeItemListEnd


%-----------CERTIFICATIONS & CREDENTIALS-----------
\section{Certifications \& Credentials}
  \resumeItemListStart
    \resumeItem{Centri Security --- Blue Team Junior Analyst (BTJA)}
    \resumeItem{Anthropic --- Model Context Protocol (MCP) Developer}
    \resumeItem{Kaggle --- Machine Learning Certification}
    \resumeItem{Cisco --- Completed specialized credentials in Application Vulnerability Exploitation, Post-Exploitation Techniques, and Wireless Security}
    \resumeItem{Deloitte --- Cybersecurity Job Simulation (Log Analysis \& Incident Response)}
    \resumeItem{Tata Group --- Identity \& Access Management (IAM) Simulation}
  \resumeItemListEnd

%-----------LANGUAGES-----------
\section{Languages}
  \resumeItemListStart
    \resumeItem{English (Professional), Telugu (Native)}
  \resumeItemListEnd

\end{document}
